\documentclass{article}

\usepackage{fancyhdr}
\usepackage{extramarks}
\usepackage{amsmath}
\usepackage{amsthm}
\usepackage{amsfonts}
\usepackage{tikz}
\usepackage[plain]{algorithm}
\usepackage{algpseudocode}
\usepackage{listings}
\usepackage{xcolor}

\usetikzlibrary{automata,positioning}

\newcommand{\overbar}[1]{\mkern 1.5mu\overline{\mkern-1.5mu#1\mkern-1.5mu}\mkern 1.5mu}

%
% Basic Document Settings
%

\topmargin=-0.45in
\evensidemargin=0in
\oddsidemargin=0in
\textwidth=6.5in
\textheight=9.0in
\headsep=0.25in

\linespread{1.1}

\pagestyle{fancy}
\lhead{Yousef Alaa Awad}
\chead{\hmwkClass\: \hmwkTitle}
\rhead{\firstxmark}
\lfoot{\lastxmark}
\cfoot{\thepage}

\renewcommand\headrulewidth{0.4pt}
\renewcommand\footrulewidth{0.4pt}

\setlength\parindent{0pt}

%
% Create Problem Sections
%

\setcounter{secnumdepth}{0}
\newcounter{partCounter}
\newcounter{homeworkProblemCounter}
\setcounter{homeworkProblemCounter}{1}

\newcommand{\hmwkTitle}{Homework\ \#3}
\newcommand{\hmwkDueDate}{February 20, 2026}
\newcommand{\hmwkClass}{Advanced Computer Architecture}

%
% Title Page
%

\title{
    \vspace{2in}
    \textmd{\textbf{\hmwkClass:\ \hmwkTitle}}\\
    \normalsize\vspace{0.1in}
}

\author{Yousef Alaa Awad}

% --- Define Custom Colors ---
\definecolor{mipsbg}{HTML}{F8F9FA}         % Light gray background
\definecolor{mipscomment}{HTML}{228B22}    % Forest Green for comments
\definecolor{mipsinstr}{HTML}{0055CC}      % Deep Blue for instructions
\definecolor{mipsdirective}{HTML}{800080}  % Purple for directives
\definecolor{mipsregister}{HTML}{D2691E}   % Orange/Chocolate for registers
\definecolor{mipsstring}{HTML}{B22222}     % Firebrick Red for strings
\definecolor{linegray}{HTML}{A9A9A9}       % Gray for line numbers

% --- Define MIPS Language ---
\lstdefinelanguage{MIPS}{
  % Class 1: Instructions (Standard and Pseudo)
  morekeywords=[1]{add, addi, addu, sub, subu, and, andi, or, ori, xor, xori, nor, sll, srl, sra, sllv, srlv, srav, slt, sltu, slti, sltiu, j, jr, jal, beq, bne, blez, bgtz, bltz, bgez, lw, sw, lb, sb, lh, sh, lui, mfhi, mflo, mult, div, syscall, la, li, move, b, bal, not, nop},
  % Class 2: Assembler Directives
  morekeywords=[2]{.align, .ascii, .asciiz, .byte, .data, .double, .extern, .float, .globl, .half, .kdata, .ktext, .space, .text, .word},
  % Class 3: Registers
  morekeywords=[3]{\$zero, \$at, \$v0, \$v1, \$a0, \$a1, \$a2, \$a3, \$t0, \$t1, \$t2, \$t3, \$t4, \$t5, \$t6, \$t7, \$s0, \$s1, \$s2, \$s3, \$s4, \$s5, \$s6, \$s7, \$t8, \$t9, \$k0, \$k1, \$gp, \$sp, \$fp, \$ra},
  % Comments and Strings
  morecomment=[l]{\#},
  morestring=[b]",
  sensitive=false, % MIPS assembly is generally case-insensitive for instructions
}

% --- Define MIPS Style ---
\lstdefinestyle{mipsStyle}{
  language=MIPS,
  backgroundcolor=\color{mipsbg},
  basicstyle=\ttfamily\small,
  keywordstyle=[1]\color{mipsinstr}\bfseries,     % Instructions
  keywordstyle=[2]\color{mipsdirective}\bfseries, % Directives
  keywordstyle=[3]\color{mipsregister},           % Registers
  commentstyle=\color{mipscomment}\itshape,       % Comments
  stringstyle=\color{mipsstring},                 % Strings
  numbers=left,
  numberstyle=\tiny\color{linegray},
  stepnumber=1,
  numbersep=10pt,
  tabsize=4,
  showspaces=false,
  showstringspaces=false,
  breaklines=true,
  frame=single,
  rulecolor=\color{lightgray},
  captionpos=b,
}

% Problems start here
\begin{document}

\maketitle
\pagebreak

\section{1: ISA}
Suppose we have a function that can be described in MIPS assembly as below. Assume that \$a0 is used for the input and initially contains 20, a positive integer. \$v0 is used for returning the output value.
\begin{center}
	\begin{minipage}{0.3\textwidth}
		\begin{lstlisting}[style=mipsStyle]
begin:
  addi $t0, $zero, 1
  addi $t1, $zero, 1
loop:
  slt $t2, $a0, $t1
  bne $t2, $zero, finish
  add $t0, $t0, $t1
  addi $t1, $t1, 2
  j loop
finish:
  add $v0, $t0, $zero
\end{lstlisting}
	\end{minipage}
\end{center}

\subsection{Assume we just finished the above program. What are the values in the registers?}
For this one, I'll be doing a loop table for all 10 loops that will occur until \$a0 is finally less than \$t1 and the loop breaks.
\begin{table}[h!]
	\centering
	\renewcommand{\arraystretch}{1.5} % Adds a bit of padding for readability
	\begin{tabular}{|c|c|c|c|c|c|}
		\hline
		\textbf{Loop Iteration} & \textbf{\$t0} & \textbf{\$t1} & \textbf{\$t2} & \textbf{\$v0} & \textbf{\$a0} \\
		\hline
		Initial                 & 1             & 1             & 0             & 0             & 20            \\
		\hline
		1                       & 2             & 3             & 0             & 0             & 20            \\
		\hline
		2                       & 5             & 5             & 0             & 0             & 20            \\
		\hline
		3                       & 10            & 7             & 0             & 0             & 20            \\
		\hline
		4                       & 17            & 9             & 0             & 0             & 20            \\
		\hline
		5                       & 26            & 11            & 0             & 0             & 20            \\
		\hline
		6                       & 37            & 13            & 0             & 0             & 20            \\
		\hline
		7                       & 50            & 15            & 0             & 0             & 20            \\
		\hline
		8                       & 65            & 17            & 0             & 0             & 20            \\
		\hline
		9                       & 82            & 19            & 0             & 0             & 20            \\
		\hline
		\textbf{10}             & \textbf{101}  & \textbf{21}   & \textbf{1}    & \textbf{101}  & \textbf{20}   \\
		\hline
	\end{tabular}
	\caption{Register trace across 10 loops.}
\end{table}
\newline
Therefore the final loop marker of 10 is the final values of the registers.

\pagebreak
\section{2: ISA}
Suppose that we would like to compute:
$$ C = A + B $$
$$ A = C - B $$

\subsection{Assume that all of A, B, C, and D are initially in memory. What is the correct sequence of instructions in the Stack ISA?}
It would look like the following:
\begin{center}
	\begin{minipage}{0.7\textwidth}
		\begin{lstlisting}
PUSH A # Push A into stack
PUSH B # Push B ontop of stack
ADD    # add the top 2 stack items (A, B)
POP C  # Pop the top item of the stack (A + B) into C
PUSH C # Push C into Stack
PUSH B # Push B onto stack
SUB    # Subtract B from C 
POP A  # Pop from stack into A
\end{lstlisting}
	\end{minipage}
\end{center}

\end{document}
